\section{Related Work}
\label{sec:related work}
\textbf{Any-Order Density Estimation}
Neural Autoregressive Distribution Estimation (NADE)~\cite{uria2016neural} initially leveraged implicit position awareness within its MLP architecture, parameterizing conditionals $p(\boldsymbol{x}_{o_d} | p(\boldsymbol{x}_{o<d})$ for an arbitrary ordering $o$ using a weight-sharing scheme inspired by RBMs. Masked Autoencoder for Distribution Estimation (MADE)~\cite{germain2015made} adapted standard autoencoders by carefully masking connections to enforce autoregressive properties for arbitrary orders. Autoregressive Diffusion Models (ARDMs)~\cite{hoogeboom2021autoregressive} integrated principles from diffusion processes, training a shared network on masked modeling objective. Arbitrary Conditional Distributions with Energy (ACE)~\cite{strauss2021arbitrary} utilizes energy-based models to estimate arbitrary conditionals $p(\boldsymbol{x}_i | \boldsymbol{x}_S)$. Training and inference on any-order autoregressive models the right way~\cite{shih2022training} addresses model redundancy and training inefficiency. They propose training on a smaller set of univariate conditionals $p(\boldsymbol{x}_i | \boldsymbol{x}_S)$ and upweighting the training loss for conditionals expected to be frequent during inference, leading to improved likelihoods without sacrificing tractable inference for arbitrary conditional queries.  InDIGO~\cite{gu2019insertion} introduces a novel insertion-based decoding algorithm for Transformers, enabling flexible sequence generation in arbitrary orders. While InDIGO demonstrates adaptive generation, our work further investigates the implications of such order-agnostic capabilities within the discrete diffusion framework and its architectural consequences for likelihood modeling. Train for the Worst, Plan for the Best~\cite{kim2025train} focuses on the trade-off between MDM's training complexity and its inference flexibility, while our paper focuses on decoupling the generative formulation from the underlying architecture.

\textbf{Existing Decoder-only Any-Order Model} Adapting decoder-only Transformers for any-order tasks has led to distinct approaches. XLNet~\cite{yang2019xlnet} trained the model to predict tokens in all possible factorization orders to capture bidirectional context. It incorporates the target position $z_t$ using a two-stream self-attention mechanism (content stream $h_{z_t}$ and query stream $g_{z_t}$), a mechanism that differs significantly from mainstream decoder architectures. $\sigma$-GPT~\cite{pannatier2024sigma} enables any-order generation in a standard GPT architecture by incorporating the target position through an additional concatenated target positional encoding. Specifically, to predict token $\boldsymbol{x}_{\sigma(t+1)}$, the model receives the current token $\boldsymbol{x}_{\sigma(t)}$, its original position $\sigma(t)$, and the original position of the target token $\sigma(t+1)$, allowing any generation order.

\textbf{Randomized Image Generation} Randomized Autoregressive modeling (RAR)~\cite{yu2024randomized} followed $\sigma$-GPT \cite{pannatier2024sigma} and trained a standard autoregressive model with an additional target position encoding on permuted image tokens with annealing during training. 
Rand-AR~\cite{pang2024randar} is a decoder-only visual model trained on randomly permuted image tokens, using an explicit position instruction token (which doubles the token length) before each image token to specify the spatial location of the next token to be predicted.
While RAR and Rand-AR both explored training on permuted sequences, they primarily focused on the visual domain and differ from our work in several key aspects.
\textbf{(1)} these methods operate exclusively on image tokens. The statistical properties and inherent structures of image tokens (e.g., local spatial correlations) are substantially different from those of language tokens, which exhibit more complex, long-range semantic and grammatical dependencies.
\textbf{(2)} their approach to generation and the theoretical framework varies. RAR, despite training with permutations, ultimately anneals towards and generates images using a conventional raster scan order. Rand-AR does consider parallel decoding by predicting tokens for specified positions; however, it does not explicitly investigate or establish the connection between its order-agnostic generation and the principles of discrete diffusion models, a central aspect of our study.
\textbf{(3)} the evaluation focus of these vision-centric works is predominantly on generative quality metrics such as Fréchet Inception Distance (FID) and Inception Score (IS). In contrast, our research places a strong emphasis on understanding the fundamental differences in likelihood modeling capabilities (e.g., perplexity) between standard autoregressive and order-agnostic/masked diffusion paradigms, particularly when controlling for architectural choices.